% Esqueleto didático adaptado do pré-projeto fornecido pelo aluno.
\section{Procedimentos metodológicos}

Quanto à abordagem, a pesquisa é predominantemente quantitativa porque se baseará na coleta, estruturação, integração e análise de bases de dados numéricas e categóricas relacionadas a passagens aéreas, tais como datas de solicitação, emissão, cotação e viagem, origem, destino, valores pagos, preços cotados, companhia aérea, antecedência de compra, características da rota e atributos temporais. A abordagem quantitativa é adequada porque permite mensurar padrões, comparar distribuições, identificar associações entre variáveis, avaliar desempenho de modelos preditivos e testar hipóteses relacionadas à recomendação de janelas de compra. Conforme Creswell e Creswell (2021), a pesquisa quantitativa é apropriada quando se pretende examinar relações entre variáveis mensuráveis e analisar dados por meio de procedimentos estatísticos.

Será calculada a antecedência de compra, em dias, considerando a diferença entre a data de emissão ou compra e a data da viagem. Quando disponíveis, também poderão ser calculadas outras medidas de antecedência, como antecedência da solicitação, antecedência da autorização e tempo decorrido entre solicitação e emissão. Essas variáveis serão importantes para distinguir restrições administrativas internas de comportamento tarifário externo.

A Figura~\ref{fig:fluxo-fontes} representa as três fontes planejadas. A Tabela~\ref{tab:fontes-planejadas} distingue o que cada uma poderá informar; ela não contém dados empíricos coletados.

\begin{figure}[htb]
\caption{Fontes previstas para a abordagem analítica}
\label{fig:fluxo-fontes}
\centering
\fbox{\parbox{0.8\textwidth}{Registros de compras realizadas + dados públicos + cotações prospectivas $\longrightarrow$ base analítica $\longrightarrow$ avaliação de padrões e incerteza.}}
\fonte{Elaboração do autor a partir do pré-projeto (2026).}
\end{figure}

\begin{table}[htb]
\caption{Papel analítico das fontes planejadas}
\label{tab:fontes-planejadas}
\centering
\begin{tabular}{p{0.3\textwidth}p{0.55\textwidth}}
\hline
Fonte & Papel previsto \\
\hline
Compras realizadas & Caracterizar gasto observado e processo institucional. \\
Dados públicos & Contextualizar padrões agregados do mercado. \\
Cotações prospectivas & Observar preços repetidos sob condições controladas. \\
\hline
\end{tabular}
\fonte{Elaboração do autor a partir do pré-projeto (2026).}
\end{table}
